% ══════════════════════════════════════════════════════════════════════════
% §1  The Concrete Problem: The 1-D Heat Rod
% ══════════════════════════════════════════════════════════════════════════
\section{The Concrete Problem: The 1-D Heat Rod}

\subsection{The Physics}

We have a metal rod of length $L$.  Both ends are held at $0^\circ$C.  A heat
source $f(x)$ is applied along the rod.  We want the \textbf{steady-state}
temperature $u(x)$, which satisfies the Poisson equation
\[
    -u''(x) = f(x), \qquad u(0) = u(L) = 0.
\]
Discretise the rod into $n$ interior points with spacing $h = L/(n+1)$.
Replacing $u''$ with the second-order central difference gives
\[
    \frac{-u_{i-1} + 2u_i - u_{i+1}}{h^2} = f_i,
    \qquad i = 1,\dots,n.
\]
This is a linear system $Au = h^2 f$ with the \textbf{tridiagonal} matrix
\begin{equation}\label{eq:tridiag}
    A = \begin{bmatrix}
         2 & -1 &        &        \\
        -1 &  2 & -1     &        \\
           & \ddots & \ddots & \ddots \\
           &        & -1     &  2
    \end{bmatrix}
    \in \mathbb{R}^{n\times n}.
\end{equation}

% ── Case A ────────────────────────────────────────────────────────────────
\subsection{Case A: The $3\times 3$ Problem (By Hand)}

Set $n=3$ and $f_i=1$ for all~$i$.  Then $b = h^2[1,1,1]^T$ and
\[
    A = \begin{bmatrix} 2 & -1 & 0 \\ -1 & 2 & -1 \\ 0 & -1 & 2 \end{bmatrix}.
\]
You can solve this by Gaussian elimination in a few lines.
The exact solution is $u = \tfrac{1}{4}[1, \tfrac{3}{2}, 1]^T$ (after scaling
by $h^2$).  Computing $A^{-1}$ directly already requires tracking nine entries
and two row-reduction passes---even a single rounding mistake
in one pivot cascades through every subsequent row.

\begin{remark}
For a \emph{tridiagonal} system, Gaussian elimination is $O(n)$.
General $LU$ is $O(n^3)$.  Structure \emph{always} matters.
\end{remark}

The following C++ snippet solves Case~A with Eigen:

\lstinputlisting[caption={Case A: $3\times 3$ heat rod (dense).},
                 label={lst:heat3}]
                {code/heat_rod_3x3.cpp}

% ── Case B ────────────────────────────────────────────────────────────────
\subsection{Case B: The $1{,}000{,}000 \times 1{,}000{,}000$ Problem}

Now the rod has $n = 10^6$ interior points.

\begin{itemize}[nosep]
    \item A \textbf{dense} $10^6 \times 10^6$ matrix of doubles needs
          $10^{12} \times 8\;\text{B} = 8\;\text{TB}$ of RAM.
    \item The tridiagonal $A$ has only $\approx 3n$ non-zeros---about
          $24\;\text{MB}$.
\end{itemize}

\begin{warningbox}
If you try to allocate a dense $10^6 \times 10^6$ matrix,
your program will be killed by the OS before it prints ``Hello.''
\end{warningbox}

This single observation motivates every topic that follows:
\emph{sparse storage} (\S\ref{sec:sparse}),
\emph{iterative solvers} (\S\ref{sec:gmres}),
and \emph{hardware awareness} (\S\ref{sec:hardware}).

\lstinputlisting[caption={Case B: $10^6$ heat rod (sparse, direct solve).},
                 label={lst:heat1M}]
                {code/heat_rod_1M.cpp}
